% ============ 第5章 §5.5 Schwartz 分布的微分 (原书 pp.118-119) ============

\section{Schwartz 分布的微分}\label{sec:5.5}
\index{0symbolSdiff@$S'(\mathbb{R}^{n})$ 中的微分}




对函数可以考虑种种微分, 但对 Schwartz 分布基本上考虑偏导数. 交换偏导数的顺序需要某种条件, 但重要的是在 Schwartz 分布的情形可以无条件地进行. 与函数的情形一样, 通过平移取差分再取极限可以得到偏导数. 但是, 这样必须讨论极限的存在性, 而且以差分定义的分布本身在很多情形下找不到那么大的用处, 所以基于 $\mathcal{S}(\mathbb{R}^{n})$ 通过积分
\[
F_{f}: g \in \mathcal{S}(\mathbb{R}^{n}) \mapsto \int_{\mathbb{R}^{n}}f(x)g(x)\dd x \quad (f \in \mathcal{S}(\mathbb{R}^{n}))
\]
可以嵌入 $\mathcal{S}'(\mathbb{R}^{n})$ 这一点, 采用用分部积分定义的方法. 如例~\ref{exa:5.4} 所示, 若 $f, g \in \mathcal{S}(\mathbb{R}^{n})$, 则由分部积分,
\[
\int_{\mathbb{R}^{n}}\partial^{\alpha}f(x) \cdot g(x)\dd x = (-1)^{|\alpha|}\int_{\mathbb{R}^{n}}f(x) \cdot \partial^{\alpha}g(x)\dd x
\]
成立, 这一点是关键.

想要对 $f \in \mathcal{S}'(\mathbb{R}^{n})$ 定义 $\partial^{\alpha}f \in \mathcal{S}'(\mathbb{R}^{n})$, 但由于 $\mathcal{S}'(\mathbb{R}^{n})$ 是从 $\mathcal{S}(\mathbb{R}^{n})$ 到 $\mathbb{C}$ 的映射, 必须考虑 $\partial^{\alpha}f$ 在 $\varphi$ 处的值. 基于这一点, 做下面的定义.

\begin{definition}[$\mathcal{S}'(\mathbb{R}^{n})$ 中的微分]\label{def:5.6}
对 $f \in \mathcal{S}'(\mathbb{R}^{n})$ 与多重指标 $\alpha \in \mathbb{N}_{0}^{n}$, 定义
\[
\langle\partial^{\alpha}f, \varphi\rangle \equiv (-1)^{|\alpha|}\langle f, \partial^{\alpha}\varphi\rangle.
\]
\end{definition}

Schwartz 分布的平移也一样, 由于要对 $f \in \mathcal{S}'(\mathbb{R}^{n})$ 定义 $f(\cdot + h) \in \mathcal{S}'(\mathbb{R}^{n})$, 必须考虑 $f(\cdot + h)$ 在 $\varphi$ 处的值.
\[
\langle f(\cdot + h), \varphi\rangle = \langle f, \varphi(\cdot - h)\rangle \quad (\varphi \in \mathcal{S}(\mathbb{R}^{n}))
\]
这样定义.

分布的微分与通常的微分有很大不同. 实际上, 如下面的定理所示, 在 $\mathcal{S}'(\mathbb{R}^{n})$ 中可以自由交换极限与微分的顺序.

\begin{theorem}[$\mathcal{S}'(\mathbb{R}^{n})$ 中极限与微分顺序的交换]\label{thm:5.10}
设 $\alpha, \beta \in \mathbb{N}_{0}^{n}$, $f \in \mathcal{S}'(\mathbb{R}^{n})$.
(1) 若 $\mathcal{S}'(\mathbb{R}^{n})$ 的点列 $\{f_{j}\}_{j=1}^{\infty}$ 收敛于 f, 则 $\lim_{j \to \infty}\partial^{\alpha}f_{j} = \partial^{\alpha}f$ 成立.
(2) $\partial^{\alpha}(\partial^{\beta}f) = \partial^{\beta}(\partial^{\alpha}f) = \partial^{\alpha+\beta}f$ 成立.
\end{theorem}

\begin{proof}
(1) 当 $j \to \infty$ 时 $f_{j} \to f$ 在 $\mathcal{S}'(\mathbb{R}^{n})$ 的意义下成立, 与对任意 $\varphi \in \mathcal{S}(\mathbb{R}^{n})$, $\langle f_{j}, \varphi\rangle \to \langle f, \varphi\rangle$ 成立是等价的. 因此, 只需证明 $\langle\partial^{\alpha}f_{j}, \varphi\rangle \to \langle\partial^{\alpha}f, \varphi\rangle$, 而这与 $(-1)^{|\alpha|}\langle f_{j}, \partial^{\alpha}\varphi\rangle \to (-1)^{|\alpha|}\langle f, \partial^{\alpha}\varphi\rangle$ 同义.

(2) 看 $\varphi \in \mathcal{S}(\mathbb{R}^{n})$ 处的值. 最左边是
\begin{align*}
\langle\partial^{\alpha}(\partial^{\beta}f), \varphi\rangle &= (-1)^{|\alpha|}\langle\partial^{\beta}f, \partial^{\alpha}\varphi\rangle \\
&= (-1)^{|\alpha+\beta|}\langle f, \partial^{\beta}\partial^{\alpha}\varphi\rangle \\
&= (-1)^{|\alpha+\beta|}\langle f, \partial^{\alpha+\beta}\varphi\rangle.
\end{align*}
$\langle\partial^{\beta}(\partial^{\alpha}f), \varphi\rangle$ 也变形为同样的式子, 所以得到结论.
\end{proof}

分布的微分的定义可以如下改写.

\begin{theorem}\label{thm:5.11}
对 $f \in \mathcal{S}'(\mathbb{R}^{n})$ 的 $\partial_{x_{j}}f$ 是古典定义的自然扩张. 即, 设 $e_{j}$ 为第 j 个基本向量,
\[
\partial_{x_{j}}f = \lim_{h \to 0}\frac{1}{h}(f(\cdot + h\mathbf{e}_{j}) - f)
\]
在 $\mathcal{S}'(\mathbb{R}^{n})$ 的拓扑中成立.
\end{theorem}

\begin{proof}
对 $\varphi \in \mathcal{S}(\mathbb{R}^{n})$, 由例~\ref{exa:5.10},
\begin{align*}
\left\langle\frac{1}{h}(f(\cdot + h\mathbf{e}_{j}) - f), \varphi\right\rangle &= \left\langle f, \frac{1}{h}(\varphi(\cdot - h\mathbf{e}_{j}) - \varphi)\right\rangle \\
&\to -\langle f, \partial_{x_{j}}\varphi\rangle = \langle\partial_{x_{j}}f, \varphi\rangle
\end{align*}
成立, 所以得到结论.
\end{proof}

\begin{exercise}\label{exer:5.11}
证明 $f(t) = e^{t}\sin(e^{t})$ 可以看作 $\mathcal{S}'(\mathbb{R})$ 的元. 【提示】找出满足
\[
\left|\int_{\mathbb{R}}f(t)\varphi(t)\dd t\right| \leq Cp_{2,1}(\varphi) \quad (\varphi \in C_{\mathrm{c}}^{\infty}(\mathbb{R}))
\]
的 $C > 0$.
\end{exercise}
